% GENERATED FILE. Edit canonical lessons, then run npm run generate:books.
% canonical-source-hash: fnv1a64:e364219fcce0b7b2
% canonical-lessons: GE-C433-orientieren, GE-C433-kompass, GE-C433-himmelsrichtung, GE-C433-zurechtfinden, GE-C433-verirren

\chapter{Wo ist Norden}
\label{ch:ge-wo-ist-norden}
\hlchaptermodality{433}

\begin{chapteropening}
\textbf{By the end of this chapter:} \emph{I can find my bearings: getting my bearings, the compass, the points of the compass, finding my way and getting lost.}

Complete the last lesson of chapter 433: I can find my bearings: getting my bearings, the compass, the points of the compass, finding my way and getting lost.
\end{chapteropening}

\section[orientieren]{orientieren — to get one's bearings (mostly reflexive: sich orientieren)}

\label{lesson:GE-C433-orientieren}



\begin{quote}
Before the new one: say the German for to compensate, then the German for the compensation.
\end{quote}

\subsection*{You'll want to know: orientieren}
\textbf{orientieren} — \textquotedblleft{}to get one's bearings (mostly reflexive: sich orientieren)\textquotedblright{}.

\begin{quote}
In der fremden Stadt orientiert er sich am Fluss.
\end{quote}

\begin{quote}
\emph{In the unfamiliar city he gets his bearings from the river.}
\end{quote}

Mostly reflexive: \emph{ich orientiere mich}. Past participle without \emph{ge-}: \emph{ich habe mich orientiert}.

\begin{cousinweb}
\emph{orientieren} was borrowed from French \emph{orienter} 'to set facing east', built on Latin \emph{oriēns} 'rising (sun), the east', from \emph{orīrī} 'to rise' (Lewis \& Short). EtymDB does not list the German verb.
\end{cousinweb}

\begin{grammarlens}[title={What you've built}]
One more word for telling, explaining and describing at B1.
\end{grammarlens}

\subsection*{Your turn}
\begin{itemize}
\raggedright
  \item \emph{Say it:} \emph{orientieren}
  \item \emph{Say it:} \emph{orientieren}, once more
  \item \emph{Say it:} a sentence of your own with \emph{orientieren}
  \item \emph{Recall:} say the German for to compensate, then the German for the compensation, then say \emph{orientieren} again
\end{itemize}

\begin{tcolorbox}[breakable,colback=teal!4,colframe=teal!35!black,title={Before you move on}]
What does orientieren mean? (\textquotedblleft{}To get one's bearings (mostly reflexive: sich orientieren)\textquotedblright{}.) Say it once more.
\end{tcolorbox}
\section[der Kompass]{der Kompass — the compass}

\label{lesson:GE-C433-kompass}



\begin{quote}
Before the new one: say the German for the compensation, then the German for to get one's bearings.
\end{quote}

\subsection*{You'll want to know: der Kompass}
\textbf{der Kompass} — \textquotedblleft{}the compass\textquotedblright{}.

\begin{quote}
Im Wald hilft ihr der Kompass.
\end{quote}

\begin{quote}
\emph{In the forest the compass helps her.}
\end{quote}

Plural: \emph{die Kompasse}.

\begin{cousinweb}
\emph{Kompass} was borrowed from Italian \emph{compasso} 'compass, pair of compasses', usually traced to a Vulgar Latin verb built on Latin \emph{passus} 'step'. EtymDB lists the German word without a source.
\end{cousinweb}

\begin{grammarlens}[title={What you've built}]
One more word for telling, explaining and describing at B1.
\end{grammarlens}

\subsection*{Your turn}
\begin{itemize}
\raggedright
  \item \emph{Say it:} \emph{der Kompass}
  \item \emph{Say it:} \emph{der Kompass}, once more
  \item \emph{Say it:} a sentence of your own with \emph{der Kompass}
  \item \emph{Recall:} say the German for the compensation, then the German for to get one's bearings, then say \emph{der Kompass} again
\end{itemize}

\begin{tcolorbox}[breakable,colback=teal!4,colframe=teal!35!black,title={Before you move on}]
What does der Kompass mean? (\textquotedblleft{}The compass\textquotedblright{}.) Say it once more.
\end{tcolorbox}
\section[die Himmelsrichtung]{die Himmelsrichtung — the point of the compass, the direction}

\label{lesson:GE-C433-himmelsrichtung}



\begin{quote}
Before the new one: say the German for to get one's bearings, then the German for the compass.
\end{quote}

\subsection*{You'll want to know: die Himmelsrichtung}
\textbf{die Himmelsrichtung} — \textquotedblleft{}the point of the compass, the direction\textquotedblright{}.

\begin{quote}
Norden, Süden, Osten und Westen sind die vier Himmelsrichtungen.
\end{quote}

\begin{quote}
\emph{North, south, east and west are the four points of the compass.}
\end{quote}

Plural: \emph{die Himmelsrichtungen}.

\begin{cousinweb}
\emph{Himmelsrichtung} is built from \emph{Himmel} 'sky, heaven' and \emph{Richtung} 'direction', from \emph{richten} 'to direct'. Both are inherited, and \emph{Richtung} shares its root with \emph{recht} 'right', Proto-Germanic \emph{rehtaz}.
\end{cousinweb}

\begin{grammarlens}[title={What you've built}]
One more word for telling, explaining and describing at B1.
\end{grammarlens}

\subsection*{Your turn}
\begin{itemize}
\raggedright
  \item \emph{Say it:} \emph{die Himmelsrichtung}
  \item \emph{Say it:} \emph{die Himmelsrichtung}, once more
  \item \emph{Say it:} a sentence of your own with \emph{die Himmelsrichtung}
  \item \emph{Recall:} say the German for to get one's bearings, then the German for the compass, then say \emph{die Himmelsrichtung} again
\end{itemize}

\begin{tcolorbox}[breakable,colback=teal!4,colframe=teal!35!black,title={Before you move on}]
What does die Himmelsrichtung mean? (\textquotedblleft{}The point of the compass, the direction\textquotedblright{}.) Say it once more.
\end{tcolorbox}
\section[zurechtfinden]{zurechtfinden — to find one's way around (mostly reflexive: sich zurechtfinden)}

\label{lesson:GE-C433-zurechtfinden}



\begin{quote}
Before the new one: say the German for the compass, then the German for the point of the compass.
\end{quote}

\subsection*{You'll want to know: zurechtfinden}
\textbf{zurechtfinden} — \textquotedblleft{}to find one's way around (mostly reflexive: sich zurechtfinden)\textquotedblright{}.

\begin{quote}
Nach einer Woche findet sie sich in der Stadt gut zurecht.
\end{quote}

\begin{quote}
\emph{After a week she finds her way around the town well.}
\end{quote}

Reflexive and separable: \emph{ich finde mich zurecht}, \emph{ich habe mich zurechtgefunden}.

\begin{cousinweb}
\emph{zurechtfinden} is built from \emph{zurecht} 'into the right place', from \emph{recht} 'right' (Proto-Germanic \emph{rehtaz}), and \emph{finden} 'to find'; both are inherited.
\end{cousinweb}

\begin{grammarlens}[title={What you've built}]
One more word for telling, explaining and describing at B1.
\end{grammarlens}

\subsection*{Your turn}
\begin{itemize}
\raggedright
  \item \emph{Say it:} \emph{zurechtfinden}
  \item \emph{Say it:} \emph{zurechtfinden}, once more
  \item \emph{Say it:} a sentence of your own with \emph{zurechtfinden}
  \item \emph{Recall:} say the German for the compass, then the German for the point of the compass, then say \emph{zurechtfinden} again
\end{itemize}

\begin{tcolorbox}[breakable,colback=teal!4,colframe=teal!35!black,title={Before you move on}]
What does zurechtfinden mean? (\textquotedblleft{}To find one's way around (mostly reflexive: sich zurechtfinden)\textquotedblright{}.) Say it once more.
\end{tcolorbox}
\section[verirren]{verirren — to get lost (mostly reflexive: sich verirren)}

\label{lesson:GE-C433-verirren}



\begin{quote}
Before the new one: say the German for the point of the compass, then the German for to find one's way around.
\end{quote}

\subsection*{You'll want to know: verirren}
\textbf{verirren} — \textquotedblleft{}to get lost (mostly reflexive: sich verirren)\textquotedblright{}.

\begin{quote}
Im Nebel haben sich die Wanderer verirrt.
\end{quote}

\begin{quote}
\emph{In the fog the hikers got lost.}
\end{quote}

Reflexive: \emph{ich habe mich verirrt}. On foot, people also say \emph{sich verlaufen}.

\begin{cousinweb}
\emph{verirren} is built from \emph{ver-} and \emph{irren} 'to wander, to go astray', inherited from Old High German \emph{irren} and Proto-Germanic \emph{irzijaną} (EtymDB).
\end{cousinweb}

\begin{grammarlens}[title={What you've built}]
One more word for telling, explaining and describing at B1.
\end{grammarlens}

\subsection*{Your turn}
\begin{itemize}
\raggedright
  \item \emph{Say it:} \emph{verirren}
  \item \emph{Say it:} \emph{verirren}, once more
  \item \emph{Say it:} a sentence of your own with \emph{verirren}
  \item \emph{Recall:} say the German for the point of the compass, then the German for to find one's way around, then say \emph{verirren} again
\end{itemize}

\begin{tcolorbox}[breakable,colback=teal!4,colframe=teal!35!black,title={Before you move on}]
What does verirren mean? (\textquotedblleft{}To get lost (mostly reflexive: sich verirren)\textquotedblright{}.) Say it once more.
\end{tcolorbox}
